\section{Employment}

\begin{frame}{Employment and Transfers}
    \mycomment{
    The population consists of those who are earning a living, and those who do not. If anyone who can work and wants to work is employed, we speak of full employment.

    The diagram does not show the correct proportions, the sizes of the boxes are made to fit the text. 
    
    The actual unemployment rate is typically below 5\%, thus this diagram overrepresents that block!
    }
    \showmytikz[]{../shared/tikz/tikz-redistribution-realistic}
\end{frame}

\begin{frame}{Employment and Transfers}
    \mycomment{
    There are, however, always some people between jobs, or otherwise unemployed while job seeking.
    
    The exact duration of such job seeking is also a matching problem. The more specialised the division of labour in a society is, the harder it is to find exact matching jobs. Paradoxically, the more skilled a worker is, the longer the search might take, or it might take re-skilling.

    Transitions, such as closure of an industry or shifts from one to the other, often result in frictions on the labour markets.
    }
    \showmytikz[1]{../shared/tikz/tikz-redistribution-realistic}
\end{frame}

\begin{frame}{Employment and Transfers}
    \mycomment{
    Only those employed have an income from labour. However, eveyone was young once and will be old one day; everyone employed now risks unemployment or illness.
    }
    \showmytikz[1,2]{../shared/tikz/tikz-redistribution-realistic}
\end{frame}

\begin{frame}{Employment and Transfers}
    \mycomment{
    It would be desirable to allow an income to those groups of the population who don't have labour income. But where can this income come from?
    }
    \showmytikz[1,2,3]{../shared/tikz/tikz-redistribution-realistic}
\end{frame}

\begin{frame}{Employment and Transfers}
    \mycomment{
    Those who have an income must transfer a part of this to those in needs. This can be done via various channels or mechanisms. 
    
    One is insurance, such as unemployment insurance, health insurance or private pension schemes. 

    Another is taxation, in which case the labour income is payed to the state, who then provides for those in need.

    Another is privately, voluntarily or obligatory, such as parents being obliged to finance their children, or relatives to take care for their family.
    }
    \showmytikz[1,2,3,4]{../shared/tikz/tikz-redistribution-realistic}
\end{frame}

\begin{frame}{Employment and Transfers}
    \mycomment{
    It must be mentioned that there is not only labour income, but also capital income, which can be redistributed as well. Capital income is more anonymous, thus would not work as well for direct insurance redistribution.

    Some pension schemes consist of building up capital ownership rather than of direct redistribution.
    }
    \showmytikz[1,2,3,4,5,6,7]{../shared/tikz/tikz-redistribution-realistic}
\end{frame}

\begin{frame}{Employment and Transfers}
    \mytext{
    \begin{itemize}
        \item The sum of all transfers received must equal the sum of transfers paid (balance mechanic)
        \item otherwise, debt builds up
        \item this is independent of the actual mechanism of redistribution
        \item there are also issues of redistribution within earners of income (high income vs. low income)
        \item the (possible) contribution of capital gains is often overlooked
        \item this only looks at income and consumption (flows), the buildup of wealth (stocks) is out of the picture
    \end{itemize}
    }
\end{frame}



\begin{frame}{Employment}
    \showmytikz{../shared/tikz/tikz-laboursupply}
\end{frame}

\begin{frame}{Employment}
    \showmytikz{../shared/tikz/tikz-labouremployment}
\end{frame}

\begin{frame}{Employment}
    \showmytikz[1]{../shared/tikz/tikz-labouremployment}
\end{frame}

\begin{frame}{Employment}
    \showmytikz[1,2]{../shared/tikz/tikz-labouremployment}
\end{frame}

\begin{frame}{Employment}
    \showmytikz[1,3]{../shared/tikz/tikz-labouremployment}
\end{frame}

\begin{frame}{Employment}
    \showmytikz{../shared/tikz/tikz-beveridge}
    \mycomment{
    Empirically, the unemployment rate and the number of open positions, as measured by job offerings, relate to each other somewhere near a curve named `Beveridge curve'. 

    This relationship can tell whether the job market is `tight' or not.
    }
\end{frame}

\begin{frame}{Employment}
    \showmytikz[1]{../shared/tikz/tikz-beveridge}
    \mycomment{
        Firms keep posting vacancies as long as the cost to do so equals the marginal revenue of doing so; which is why it is a straight line.
    }   
\end{frame}


%% okun's law, structural unemployment, temporal unemployment, phillips curve, insider-outsider


